\documentclass[11pt]{article}
\usepackage{amsmath,amssymb,amsthm}

\providecommand{\uses}[1]{}
\providecommand{\lean}[1]{}
\providecommand{\leanok}{}

\newtheorem{theorem}{Theorem}
\newtheorem{lemma}[theorem]{Lemma}
\theoremstyle{definition}
\newtheorem{definition}[theorem]{Definition}

\begin{document}

%% =====================================================================
%% Definitions
%% =====================================================================

\begin{definition}[Lattice point]
\label{def:lattice-point}
\lean{Pick.LatticePoint}
A \emph{lattice point} in the plane is a point of $\mathbb{Z}^2 \subset
\mathbb{R}^2$.
\end{definition}

\begin{definition}[Simple lattice polygon]
\label{def:simple-polygon}
\lean{Pick.SimplePolygon}
\uses{def:lattice-point}
A \emph{simple lattice polygon} is a closed polygonal curve in $\mathbb{R}^2$
whose vertices are lattice points and whose edges intersect only at shared
endpoints, together with the closed bounded region it encloses. We write $P$
for both the curve and the region; context disambiguates.
\end{definition}

\begin{definition}[Interior lattice points]
\label{def:interior-points}
\lean{Pick.SimplePolygon.interiorLatticePoints}
\uses{def:simple-polygon}
For a simple lattice polygon $P$, write $i(P)$ for the number of lattice
points lying in the topological interior of $P$.
\end{definition}

\begin{definition}[Boundary lattice points]
\label{def:boundary-points}
\lean{Pick.SimplePolygon.boundaryLatticePoints}
\uses{def:simple-polygon}
For a simple lattice polygon $P$, write $b(P)$ for the number of lattice
points lying on the boundary curve of $P$ (vertices included).
\end{definition}

\begin{definition}[Polygon area]
\label{def:area}
\lean{Pick.SimplePolygon.area}
\uses{def:simple-polygon}
For a simple lattice polygon $P$, write $A(P)$ for the Euclidean area of the
closed region enclosed by $P$ (e.g.\ via the shoelace formula on the vertex
sequence).
\end{definition}

\begin{definition}[Lattice triangle]
\label{def:lattice-triangle}
\lean{Pick.LatticeTriangle}
\uses{def:lattice-point}
A \emph{lattice triangle} is the convex hull of three non-collinear lattice
points $v_1, v_2, v_3 \in \mathbb{Z}^2$.
\end{definition}

\begin{definition}[Primitive lattice triangle]
\label{def:primitive-triangle}
\lean{Pick.LatticeTriangle.IsPrimitive}
\uses{def:lattice-triangle}
A lattice triangle $\Delta$ is \emph{primitive} if the only lattice points of
$\mathbb{Z}^2$ contained in $\Delta$ (interior or boundary) are its three
vertices.
\end{definition}

\begin{definition}[Triangulation of a polygon]
\label{def:triangulation}
\lean{Pick.Triangulation}
\uses{def:simple-polygon, def:lattice-triangle}
A \emph{lattice triangulation} of a simple lattice polygon $P$ is a finite
collection $\mathcal{T} = \{\Delta_1, \dots, \Delta_T\}$ of lattice triangles
such that:
\begin{enumerate}
\item their union equals the closed region $P$;
\item their interiors are pairwise disjoint;
\item the intersection of any two distinct triangles is either empty, a
common vertex, or a common edge (the \emph{edge-to-edge} condition).
\end{enumerate}
The triangulation is called \emph{primitive} if every $\Delta_j$ is a
primitive lattice triangle, and \emph{full} if every lattice point of $P$
appears as a vertex of some $\Delta_j$.
\end{definition}

\begin{definition}[Plane graph of a triangulation]
\label{def:plane-graph}
\lean{Pick.Triangulation.planeGraph}
\uses{def:triangulation}
Given a lattice triangulation $\mathcal{T}$ of $P$, its \emph{plane graph}
$G(\mathcal{T})$ has vertex set equal to the set of triangle vertices and
edge set equal to the set of triangle sides. Its faces (in the plane) are the
$T$ open triangle interiors together with the unbounded exterior of $P$.
\end{definition}

%% =====================================================================
%% Sub-lemmas
%% =====================================================================

\begin{lemma}[Primitive lattice triangle has area $\tfrac12$]
\label{lem:primitive-triangle-area}
\lean{Pick.primitive_lattice_triangle_area}
\uses{def:primitive-triangle, def:area}
If $\Delta$ is a primitive lattice triangle, then its Euclidean area is
$\tfrac12$.
\end{lemma}

\begin{lemma}[Existence of a full primitive triangulation]
\label{lem:exists-primitive-triangulation}
\lean{Pick.exists_primitive_lattice_triangulation}
\uses{def:simple-polygon, def:triangulation, def:primitive-triangle}
Every simple lattice polygon $P$ admits a triangulation $\mathcal{T}$ that is
both \emph{full} (every lattice point of $P$ is a vertex of $\mathcal{T}$)
and \emph{primitive} (every triangle in $\mathcal{T}$ is primitive).
\end{lemma}

\begin{lemma}[Area additivity over a triangulation]
\label{lem:area-additive}
\lean{Pick.area_eq_sum_triangle_areas}
\uses{def:triangulation, def:area}
If $\mathcal{T} = \{\Delta_1, \dots, \Delta_T\}$ is any lattice triangulation
of $P$, then
\[
A(P) \;=\; \sum_{j=1}^{T} A(\Delta_j).
\]
In particular, if $\mathcal{T}$ is primitive, then $A(P) = T/2$.
\end{lemma}

\begin{lemma}[Vertex count]
\label{lem:vertex-count}
\lean{Pick.vertex_count}
\uses{def:plane-graph, def:interior-points, def:boundary-points,
lem:exists-primitive-triangulation}
For a full lattice triangulation $\mathcal{T}$ of $P$ with plane graph
$G(\mathcal{T})$,
\[
V \;=\; i(P) + b(P),
\]
where $V$ is the number of vertices of $G(\mathcal{T})$.
\end{lemma}

\begin{lemma}[Face count]
\label{lem:face-count}
\lean{Pick.face_count}
\uses{def:plane-graph, def:triangulation}
For a lattice triangulation $\mathcal{T}$ of $P$ with $T$ triangles,
\[
F \;=\; T + 1,
\]
where $F$ is the number of faces of the plane graph $G(\mathcal{T})$
(counting the unbounded exterior face).
\end{lemma}

\begin{lemma}[Boundary contributes $b$ edges]
\label{lem:boundary-edge-count}
\lean{Pick.boundary_edge_count}
\uses{def:plane-graph, def:boundary-points,
lem:exists-primitive-triangulation}
For a full lattice triangulation $\mathcal{T}$ of $P$, the boundary of $P$
contributes exactly $b(P)$ edges to the plane graph $G(\mathcal{T})$, one
between each pair of consecutive boundary lattice points.
\end{lemma}

\begin{lemma}[Edge double count]
\label{lem:edge-double-count}
\lean{Pick.edge_double_count}
\uses{def:plane-graph, def:triangulation, lem:boundary-edge-count}
For a lattice triangulation $\mathcal{T}$ of $P$ with $T$ triangles and plane
graph $G(\mathcal{T})$ having $E$ edges,
\[
3T \;=\; 2E - b(P),
\qquad\text{equivalently}\qquad
E \;=\; \frac{3T + b(P)}{2}.
\]
The proof counts triangle sides in two ways: each interior edge lies on two
triangles, each boundary edge on one.
\end{lemma}

\begin{lemma}[Euler's formula for the triangulation graph]
\label{lem:euler-formula}
\lean{Pick.euler_formula}
\uses{def:plane-graph}
For the plane graph $G(\mathcal{T})$ of a lattice triangulation of a simple
polygon,
\[
V - E + F \;=\; 2.
\]
\end{lemma}

\begin{lemma}[Triangle count]
\label{lem:triangle-count}
\lean{Pick.triangle_count}
\uses{lem:vertex-count, lem:face-count, lem:edge-double-count,
lem:euler-formula}
For a full primitive lattice triangulation $\mathcal{T}$ of $P$ with $T$
triangles,
\[
T \;=\; 2\,i(P) + b(P) - 2.
\]
This follows by substituting $V = i+b$, $F = T+1$, and
$E = (3T+b)/2$ into Euler's formula.
\end{lemma}

%% =====================================================================
%% Main theorem
%% =====================================================================

\begin{theorem}[Pick's Theorem]
\label{thm:pick}
\lean{Pick.pick}
\uses{def:simple-polygon, def:area, def:interior-points, def:boundary-points,
lem:primitive-triangle-area, lem:exists-primitive-triangulation,
lem:area-additive, lem:triangle-count}
Let $P$ be a simple polygon in the plane whose vertices all have integer
coordinates. Let $i$ be the number of lattice points in the interior of $P$,
and let $b$ be the number of lattice points on the boundary of $P$ (including
the vertices). Then the area $A$ of $P$ satisfies
\[
A \;=\; i + \frac{b}{2} - 1.
\]
\end{theorem}

\begin{proof}
Choose a full primitive lattice triangulation $\mathcal{T}$ of $P$ with $T$
triangles, which exists by Lemma~\ref{lem:exists-primitive-triangulation}.
By Lemma~\ref{lem:primitive-triangle-area} each triangle has area $\tfrac12$,
and by Lemma~\ref{lem:area-additive} the polygon area is the sum of triangle
areas, so $A = T/2$. By Lemma~\ref{lem:triangle-count},
$T = 2i + b - 2$. Therefore
\[
A \;=\; \frac{T}{2} \;=\; \frac{2i + b - 2}{2} \;=\; i + \frac{b}{2} - 1.
\qedhere
\]
\end{proof}

\end{document}
